% Generated by paper4/manuscript/build.py from paper4/paper/section_discussion.md (68dd278bac381a45): edit the Markdown, not this file.
\section{Discussion}
\textbf{Why no further model.} One could object that Graph WaveNet is one model among many, and a better architecture might extract what it did not. The analyses without training answer from the data rather than from the model. A table of class frequencies by hour of the week forecasts a segment without a sensor as well as the trained model without spatial exchange; adding the nearest instrumented neighbour's class adds nothing within 250 m, where every interval contains zero, and makes the forecast worse beyond. A learned model can only exploit a dependence that the data contain, and the dependence between a segment without a sensor and its neighbours, net of the calendar, is measured here as zero within the only radius where the labels agree. A further model would have to find a signal that neither the learned propagation nor the calibrated tables found, so we did not train one. The argument has a stated perimeter (§VI): one neighbour's calibrated class, and a learned aggregation of several neighbours over three hours, on labels in three classes.

\textbf{What the agreement of nearby labels is.} Two segments within 250 m share a class more often than chance. The natural reading is that congestion spreads along a corridor, so a sensor tells about its unmeasured neighbours. The forecasts say otherwise: nearby segments agree because both are congested at the same hours of the same days. The calendar is common knowledge, and it needs no sensor nearby.

\textbf{What placement can and cannot do.} Contiguous blocks leave most uncovered segments far from any sensor, while uniform placement brings almost all of them within 250 m from 20\% coverage in Madrid and London. Since proximity has no forecasting value beyond the calendar, placement should be chosen for the segments a sensor covers, where a measurement is needed for control, enforcement or detection, not for what the sensor might predict elsewhere.

\textbf{What to measure.} A loop detector says whether a segment is congested, not how congested: its reconstruction ranks levels poorly and separates congested cells almost perfectly, with a threshold that transfers between sensors. Its low agreement with probe speeds on arterials is about half a matter of matching a loop to a segment some 30 m away. Per lane, detection is as good as on the ring; averaged over the lanes of a section, it loses the episodes that occur on one lane only. A target defined on the section mean discards part of what the sensor measures.

\textbf{Melbourne.} Melbourne is the exception in several places: the calendar-only table is ahead of G0 at 20\% coverage (+4.11 [+1.42; +6.11]), the 0–250 m bin is empty or too thin to read, and G1 and G0 are above ln 3 at 10\% and 20\%. These are one fact, stated once: the network is sparse. Its grid cells measure 1,974 m against 515–519 m in Madrid and London, and 85–98\% of its hidden segments have no instrumented segment within 500 m. Its intervals are wide; the statements of E9, E12, E12-bis and E12-ter are ``undetermined'' only because of Melbourne; and no verdict of this paper is read on that city alone.

\textbf{Relation to the literature.} Graph models for traffic are usually evaluated on sensors whose own history is in training or in the input ~\cite{wu2019gwnet}. Kriging methods estimate unobserved locations from simultaneous observations, and are trained by masking nodes that remain in the network, so the masked locations shape the learned weights ~\cite{wu2021ignnk}. Our regime is that of an installation: a segment that has never had a sensor never enters the loss, the input or the normalisation, and the forecast is 15–120 minutes ahead. The project report that preceded this work ~\cite{motion2026uc04} found propagation gains with hidden segments that were training targets; our positive control reproduces that condition and the gain that comes with it. The negative result does not contradict those findings: it concerns a stricter regime. [ref: placement and coverage studies, to be chosen.]

